% ch02-cqrs-integrity
\section{نمط CQRS وخطّ المعالجة وسلامة البيانات}

في الفصل السابق عرفنا الطبقات الأربع، وكيف أنّ الطبقة الداخلية لا تعرف شيئًا عن
الخارجية، وكيف تتكلّم الطبقات مع العالم الخارجيّ عبر \en{ports \& adapters}. في هذا
الفصل ننزل درجة أعمق: كيف \emph{تسافر} أيّ عملية من لحظة وصولها للـ \en{API} حتى
تُحفَظ في قاعدة البيانات؟ ومَن الذي يحرس هذه الرحلة حتى لا تُكتَب بيانات ناقصة أو
مكرّرة أو متعارضة؟

الإجابة مبنيّة على فكرتين متكاملتين. الأولى تنظيميّة: نمط \en{CQRS} الذي يفصل عمليّات
\emph{الكتابة} عن عمليّات \emph{القراءة}، ويمرّر كلّ عملية عبر خطّ معالجة موحّد. والثانية
دفاعيّة: مجموعة ضوابط في طبقة البيانات تضمن أنّ الصفّ الذي وصل للقاعدة صحيح ومتّسق. لو
فهمت هذين المحورين تكون فهمت العمود الفقريّ للنظام كلّه.

% ====================================================================
\subsection{نمط CQRS: كل عملية إمّا أمر أو استعلام}

\en{CQRS} اختصار لـ \en{Command Query Responsibility Segregation}، أي «فصل مسؤوليّات
الأمر عن الاستعلام». الفكرة بسيطة جدًّا: أيّ شيء يفعله النظام يقع في واحدة من خانتين
فقط، لا ثالث لهما.

\begin{idea}[الفكرة: أمر أم استعلام؟]
\textbf{الأمر (\en{Command})} يُغيّر حالة النظام — يكتب، يعدّل، أو يحذف. أمثلة:
\code{AcceptBookingCommand} (قبول حجز)، \code{CreatePostCommand} (إنشاء منشور).
عادةً يُرجِع لا شيء أو نتيجة صغيرة.\\[4pt]
\textbf{الاستعلام (\en{Query})} يقرأ الحالة \emph{دون أيّ أثر جانبيّ} — لا يكتب حرفًا
واحدًا. أمثلة: \code{GetMyRecommendationsQuery}، \code{GetBookingRecordingsQuery}.
كل طرف من الطرفين له نوع طلب مستقلّ، ومعالِج (\en{handler}) مستقلّ، وقواعد تحقّق
مستقلّة.
\end{idea}

كل هذه الطلبات — أوامر واستعلامات — تُرسَل عبر مكتبة واحدة اسمها \en{MediatR}. أنت
تبني كائن الطلب وتناديه بسطر واحد \code{mediator.Send(request)}، وهي تتكفّل بالعثور على
المعالِج الصحيح وتشغيله. النتيجة المباشرة لهذا التصميم أنّ الـ \en{controller} يصبح
\emph{رفيعًا جدًّا}: لا منطق عمل بداخله، فقط يستقبل الطلب ويحوّله.

\textbf{هنشوف إيه:} الـ \en{endpoint} المسؤول عن قبول الحجز في \code{BookingsController}.
لاحظ أنّ كل عمله الحقيقيّ سطرٌ واحد.

\begin{english}
\begin{lstlisting}[language={[Sharp]C}]
[HttpPost("{id:guid}/accept")]
public async Task<IActionResult> AcceptBooking(
    Guid id,
    [FromBody] AcceptBookingCommand command,
    CancellationToken cancellationToken)
{
    if (id != command.BookingId)
        return BadRequest("Route booking id does not match body booking id.");

    await _sender.Send(command, cancellationToken);   // one line
    return NoContent();
}
\end{lstlisting}
\end{english}

شرح سريع سطرًا بسطر لأهمّ ما في المقتطف:
\begin{itemize}
\item السمة \code{[HttpPost(...)]} تربط المسار \en{/accept} بهذه الدالة.
\item الطلب القادم من العميل يُربَط مباشرةً بكائن \code{AcceptBookingCommand} — أي أنّ
  «شكل الطلب» و«الأمر» شيء واحد.
\item السطر \code{await \_sender.Send(command, ...)} هو كل المنطق: أرسِل الأمر، و\en{MediatR}
  ستجد المعالِج وتشغّله. الـ \en{controller} لا يعرف \emph{كيف} يُقبَل الحجز إطلاقًا.
\item يتحقّق فقط من اتّساق بسيط (تطابق المُعرّف في المسار والجسد) ثم يُعيد \en{204 NoContent}.
\end{itemize}

\begin{why}
لماذا نُبقي الـ \en{controller} رفيعًا؟ لأنّ كل منطق العمل يصير في صنف معالِج مستقلّ
قابل للاختبار وحده دون تشغيل خادم \en{HTTP}. كل ميزة جديدة تعني إضافة \emph{معالِج جديد}،
لا تعديل صنف خدمة ضخم قائم. وهذا يفتح الباب لخطّ معالجة موحّد يلفّ \emph{كل} معالِج
تلقائيًّا — وهو موضوعنا التالي.
\end{why}

% ====================================================================
\subsection{مثال حيّ: أمر قبول الحجز}

لنتتبّع أمرًا حقيقيًّا كاملًا. عندما يقبل المستشار حجزًا مدفوعًا، يجب أن تحدث عدّة أشياء
بترتيب دقيق: تحميل الحجز، سحب المبلغ المحجوز من بطاقة الطالب عبر \en{Stripe}، تجهيز غرفة
اجتماع مرئيّ عبر \en{ACS}، حفظ كل ذلك، ثم إعلان أنّ الحجز تأكّد. كل هذا يعيش في معالِج
واحد.

أوّلًا، الأمر نفسه — لاحظ أنّه \code{record} صغير جدًّا يحمل مُعرّف الحجز فقط، ومُزخرَف
بالسمة \code{[Auditable]}:

\textbf{هنشوف إيه:} تعريف \code{AcceptBookingCommand} والسمة التي تجعله قابلًا للتدقيق.

\begin{english}
\begin{lstlisting}[language={[Sharp]C}]
[Auditable(
    AuditAction.BookingAccepted,
    "ConsultantBooking",
    TargetIdProperty = nameof(BookingId),
    SummaryTemplate = "Booking accepted: {TargetId}")]
public sealed record AcceptBookingCommand(Guid BookingId) : IRequest;
\end{lstlisting}
\end{english}

الأمر مجرّد بيانات: مُعرّف حجز واحد. أمّا السمة \code{[Auditable]} فتقول للنظام «سجّل هذا
الأمر تلقائيًّا في سجلّ التدقيق باسم الفعل \en{BookingAccepted}» — وسنرى آليّتها بعد قليل.
الآن المعالِج (مختصَرًا مع الحفاظ على الخطوات الحقيقيّة):

\textbf{هنشوف إيه:} قلب \code{AcceptBookingCommandHandler}: تحميل، سحب الدفع، تجهيز الغرفة،
حفظ، ثم نشر الحدث.

\begin{english}
\begin{lstlisting}[language={[Sharp]C}]
public async Task Handle(AcceptBookingCommand request, CancellationToken ct)
{
    // 1. Load the aggregate (with its Payment)
    var booking = await _context.Bookings
        .Include(b => b.Payment)
        .FirstOrDefaultAsync(b => b.Id == request.BookingId, ct);
    // ... ownership + state guards omitted ...

    // 2. Stripe capture - take the money that was only *held* at booking time
    var capture = await _stripeService.CapturePaymentIntentAsync(
        booking.StripePaymentIntentId!, amountCents, idempotencyKey, ct);

    // 3. Provision the ACS video room for the confirmed session
    var room = await _meetingService.CreateRoomAsync(booking.Id, ct);
    booking.MeetingRoomId = room.RoomId;
    booking.Status = BookingStatus.Confirmed;

    // 4. Persist, THEN publish the domain event
    await _context.SaveChangesAsync(ct);
    await _publisher.Publish(
        new BookingConfirmedEvent(
            booking.Id, booking.StudentId, booking.ConsultantId), ct);
}
\end{lstlisting}
\end{english}

اقرأ الخطوات كقصّة قصيرة:
\begin{itemize}
\item \textbf{تحميل}: يجلب الحجز ومعه سجلّ الدفع \code{Payment} في استعلام واحد بفضل
  \code{Include}. (حُذفت هنا فحوص الملكيّة والحالة لِقِصَر المساحة، وهي موجودة في الكود
  الأصليّ.)
\item \textbf{سحب الدفع}: عند الحجز كان المبلغ \emph{محجوزًا} فقط على البطاقة لا مسحوبًا.
  الآن يُسحَب فعليًّا عبر \code{CapturePaymentIntentAsync} — سنشرح نموذج «احجز ثمّ اسحب»
  في فصل المدفوعات.
\item \textbf{تجهيز الغرفة}: يُنشئ غرفة اجتماع عبر خدمة \en{ACS} ويحفظ مُعرّفها على الحجز.
\item \textbf{حفظ ثم نشر}: يُستدعى \code{SaveChangesAsync} أوّلًا، وبعد نجاح الحفظ فقط
  يُنشَر \code{BookingConfirmedEvent}. هذا الترتيب مقصود تمامًا — نُوضّحه في قسم المعترِض.
\end{itemize}

\begin{keyp}[نقطة مهمة: المعالِج هو موطن منطق العمل]
كل القرارات — مَن يملك الحجز، أهو في الحالة الصحيحة، متى نسحب، متى ننشر الحدث — تعيش
هنا، في صنف واحد صغير قابل للاختبار. الـ \en{controller} لم يعرف شيئًا من ذلك. هذا هو
جوهر \en{CQRS} عمليًّا.
\end{keyp}

% ====================================================================
\subsection{خطّ المعالجة وترتيبه المقصود}

ذكرنا أنّ \en{MediatR} تلفّ كل معالِج. الأداة التي تفعل ذلك تُسمّى \en{pipeline behavior}:
حلقة تُغلِّف المعالِج فتنفّذ شيئًا \emph{قبله} و\emph{بعده}، تمامًا كطبقات البصلة. عندنا
أربع حلقات، وترتيب تسجيلها = ترتيب تنفيذها (الأبعد أوّلًا، الأقرب للمعالِج أخيرًا).

\textbf{هنشوف إيه:} تسجيل الحلقات الأربع في \code{DependencyInjection.cs} — الترتيب هنا هو
عين الترتيب الفعليّ عند التنفيذ.

\begin{english}
\begin{lstlisting}[language={[Sharp]C}]
services.AddMediatR(cfg =>
{
    cfg.RegisterServicesFromAssembly(assembly);
    cfg.AddOpenBehavior(typeof(LoggingBehavior<,>));      // 1 outermost
    cfg.AddOpenBehavior(typeof(ValidationBehavior<,>));   // 2
    cfg.AddOpenBehavior(typeof(PerformanceBehavior<,>));  // 3
    cfg.AddOpenBehavior(typeof(AuditBehavior<,>));        // 4 closest to handler
});
\end{lstlisting}
\end{english}

الجدول التالي يلخّص وظيفة كل حلقة وسبب موضعها:

\begin{center}
\begin{tabular}{L{3.2cm} L{4.4cm} L{5cm}}
\toprule
\textbf{\color{spPrimary}الحلقة} & \textbf{\color{spPrimary}ماذا تفعل} & \textbf{\color{spPrimary}لماذا هذا الموضع} \\
\midrule
\en{Logging} & تسجّل بداية ونهاية كل طلب مع مُعرّف المستخدم & الأبعد، لتلتقط \emph{كل} طلب حتى لو رُفِض لاحقًا \\
\en{Validation} & تشغّل كل مُتحقّقات \en{FluentValidation} وتقصر الطلب بـ 400 عند الفشل & مبكّرة عمدًا: لا نُضيّع وقتًا ولا نُدقّق طلبًا معطوبًا \\
\en{Performance} & تقيس زمن المعالِج وتُنبّه للبطيء & بعد التحقّق، لنقيس زمن العمل الحقيقيّ فقط \\
\en{Audit} & تكتب صفًّا في \code{AuditLog} للأوامر المُزخرَفة & الأقرب للمعالِج، لِتُسجّل فقط ما نُفِّذ فعلًا \\
\bottomrule
\end{tabular}
\end{center}

\begin{why}
الترتيب ليس اعتباطيًّا. \en{Logging} يجب أن يرى \emph{كل} طلب، لذا هو الأبعد.
\en{Validation} يجب أن يرفض المدخل السيّئ \emph{قبل} أن نصرف وقتًا في القياس أو نكتب صفوف
تدقيق لطلب لن يُنفَّذ أصلًا. \en{Audit} يجب أن يجلس أقرب ما يكون للمعالِج، فلا يُسجّل إلّا
فعلًا نُفِّذ حقًّا. اقلب الترتيب وستدقّق أفعالًا لم تحدث، أو تقيس زمن طلبات فاسدة.
\end{why}

ولنرَ حلقة التحقّق من الداخل — كيف «تقصر» الطلب:

\textbf{هنشوف إيه:} قلب \code{ValidationBehavior}: يجمع الأخطاء، وإن وُجدت يرمي استثناءً
قبل أن يصل الطلب للمعالِج.

\begin{english}
\begin{lstlisting}[language={[Sharp]C}]
var failures = validationResults
    .Where(r => !r.IsValid)
    .SelectMany(r => r.Errors)
    .ToList();

if (failures.Count > 0)
    throw new ValidationException(failures);   // short-circuit -> 400

return await next(cancellationToken);          // otherwise, continue inward
\end{lstlisting}
\end{english}

الفكرة أنّ \code{next(...)} هي «الحلقة التالية إلى الداخل» وصولًا للمعالِج. إن وُجد خطأ
واحد، نرمي \code{ValidationException} فورًا ولا نستدعي \code{next} أبدًا — فلا يعمل
المعالِج، ولا تُقاس مدّة، ولا يُكتَب تدقيق. هذا معنى «القصر» (\en{short-circuit}).

\begin{note}
المُتحقّقات ومُلفّات \en{AutoMapper} تُكتَشَف تلقائيًّا من التجميعة (\en{assembly}) عند
الإقلاع. أي أنّ إضافة ميزة جديدة بمُتحقّقها لا تحتاج أيّ ربط يدويّ في حاوية الاعتماديّات.
\end{note}

% ====================================================================
\subsection{التدقيق بالسمة Auditable}

الآن نفهم كيف يعمل التدقيق التلقائيّ. حلقة \en{Audit} — الأقرب للمعالِج — تنتظر انتهاء
المعالِج بنجاح، ثم تسأل: هل نوع هذا الطلب مُزخرَف بالسمة \code{[Auditable]}؟ إن كان
كذلك، تكتب صفًّا في جدول \code{AuditLog}. إن لم يكن، لا تفعل شيئًا.

\textbf{هنشوف إيه:} قلب \code{AuditBehavior}: يشغّل المعالِج، يقرأ السمة، وإن وُجدت يكتب
سجلّ التدقيق.

\begin{english}
\begin{lstlisting}[language={[Sharp]C}]
var response = await next(cancellationToken);   // run the handler first

var attr = (AuditableAttribute?)Attribute.GetCustomAttribute(
    typeof(TRequest), typeof(AuditableAttribute));
if (attr is null)
    return response;                            // no attribute -> no audit row

var targetId = ResolveTargetId(response, request, attr.TargetIdProperty);
var summary  = BuildSummary(attr.SummaryTemplate, targetId, response, request);

await audit.WriteAsync(attr.Action, attr.TargetType, targetId,
    beforeJson: null, afterJson: null, summary, cancellationToken);
return response;
\end{lstlisting}
\end{english}

نقاط جديرة بالانتباه:
\begin{itemize}
\item \code{await next(...)} أوّلًا — التدقيق يحدث \emph{بعد} نجاح المعالِج، فلا نُسجّل
  فعلًا فشل.
\item \code{GetCustomAttribute} يقرأ السمة عن نوع الطلب. لو غابت، يخرج فورًا — فالحلقة
  «لا تفعل شيئًا» لكل الأوامر غير المُعلَّمة.
\item \code{ResolveTargetId} و\code{BuildSummary} يستخرجان مُعرّف الهدف ونصًّا وصفيًّا من
  السمة (تذكّر \code{TargetIdProperty = nameof(BookingId)} في أمر قبول الحجز).
\item لو فشلت الكتابة في التدقيق، تُلتقَط ولا تُفشِل الأمر الأصليّ — فنجاح العمل لا يُلغى
  بسبب سجلّ.
\end{itemize}

\begin{keyp}[نقطة مهمة: التدقيق بالتزيين لا بالتكرار]
لتجعل أمرًا قابلًا للتدقيق تضع سطرًا واحدًا فوقه: \code{[Auditable(...)]}. لا كود مكرّر
داخل كل معالِج. عشرات الأوامر الحسّاسة في النظام (قبول حجز، موافقة على منحة، تغيير دور
مستخدم، ...) مُعلّمة هكذا، فتُسجَّل كلها بنفس الآليّة المركزيّة.
\end{keyp}

% ====================================================================
\subsection{سلامة البيانات: الحذف الناعم}

ننتقل الآن من «رحلة الطلب» إلى «حراسة البيانات». أوّل حارس هو \emph{الحذف الناعم}
(\en{soft delete}). الصفوف الحسّاسة أو ذات القيمة التاريخيّة لا تُحذَف حذفًا صلبًا أبدًا.
بدلًا من ذلك تُنفّذ الواجهة \code{ISoftDeletable} (تحمل \code{IsDeleted} و\code{DeletedAt}
و\code{DeletedByUserId})، و«الحذف» ما هو إلّا قلب \code{IsDeleted} إلى \en{true}.

لكن كيف نضمن أنّ الصفّ «المحذوف» يختفي من \emph{كل} استعلام دون أن يتذكّر كل مطوّر ذلك؟
عبر \emph{مُرشِّح استعلام عموميّ} في \en{EF Core} على كل كيان من هذا النوع:

\textbf{هنشوف إيه:} سطر المُرشِّح العموميّ — يُطبَّق على قرابة 15 كيانًا.

\begin{english}
\begin{lstlisting}[language={[Sharp]C}]
b.HasQueryFilter(a => !a.IsDeleted);   // ~15 entities carry this
\end{lstlisting}
\end{english}

\begin{idea}[الفكرة: الصفّ يختفي دون أن يُدمَّر]
بمجرّد وضع هذا المُرشِّح على الكيان، \emph{كل} استعلام \en{LINQ} يولّده \en{EF} لهذا
الكيان يُضاف إليه ضمنيًّا \code{WHERE IsDeleted = 0}. المطوّر لا يستطيع أن ينسى إضافته —
فالصفوف المحذوفة ببساطة لا تظهر في أيّ استعلام أو ربط أو عدّ، إلّا إن اختار الكود صراحةً
تجاوز المُرشِّح بـ \code{IgnoreQueryFilters()} (وهذا محصور في أدوات الإدارة و\en{GDPR}).
\end{idea}

\begin{why}
لماذا نتجنّب الحذف الصلب؟ لأنّه يكسر التاريخ المترابط (مراجعات، مدفوعات، آثار تدقيق)،
ويجعل سؤال «مَن حذف ماذا ومتى؟» مستحيلًا. الحذف الناعم يحفظ ترابط التاريخ، ويدعم مسار
تراجع/استرجاع، ومع سجلّ التدقيق يعطي تتبّعًا كاملًا. وهو أيضًا الخيار الآمن الافتراضيّ
لمنصّة ملزَمة بقواعد حفظ البيانات وحلّ النزاعات.
\end{why}

% ====================================================================
\subsection{معترِض SaveChangesAsync: مكان واحد، ضمانات كثيرة}

الحارس الثاني — والأقوى — هو تجاوُز الدالة \code{SaveChangesAsync} داخل \code{DbContext}.
هذه الدالة هي «الاختناق» الوحيد الذي يمرّ منه كل حفظ للبيانات، فجعلناها تفرض عدّة قواعد
عرضيّة (\en{cross-cutting}) دفعةً واحدة.

\textbf{هنشوف إيه:} قلب \code{SaveChangesAsync}: ختم التواريخ، نشر الأحداث بعد الحفظ،
وترجمة تعارض الفهرس الفريد.

\begin{english}
\begin{lstlisting}[language={[Sharp]C}]
foreach (var entry in ChangeTracker.Entries<AuditableEntity>())
{
    if (entry.State == EntityState.Added && entry.Entity.CreatedAt == default)
        entry.Entity.CreatedAt = auditNow;
    else if (entry.State == EntityState.Modified)
        entry.Entity.UpdatedAt = auditNow;
}

try
{
    var result = await base.SaveChangesAsync(cancellationToken);
    await DispatchDomainEventsAsync(cancellationToken);   // only after commit
    return result;
}
catch (DbUpdateException ex) when (ex.InnerException is SqlException sqlEx
       && (sqlEx.Number == 2601 || sqlEx.Number == 2627))
{
    throw new ConflictException("This action conflicts with an existing record.");
}
\end{lstlisting}
\end{english}

هذا المقتطف الصغير يفرض ثلاث ضمانات (والرابعة نشرحها بعده نصًّا):
\begin{itemize}
\item \textbf{ختم التواريخ}: قبل الحفظ يمرّ على الكيانات المتغيّرة، فيضع \code{CreatedAt}
  على الصفوف الجديدة (فقط إن لم تُضبَط، حتى تحتفظ بيانات البذر بتواريخها المقصودة)،
  و\code{UpdatedAt} على كل تعديل. بدون هذا تُحفَظ صفوف بتاريخ صفريّ يكسر العرض والترتيب.
\item \textbf{أحداث الدومين بعد الحفظ فقط}: يُستدعى \code{DispatchDomainEventsAsync}
  \emph{بعد} أن ينجح \code{base.SaveChangesAsync}. لماذا؟ حتى لا يُطلَق أثر جانبيّ
  (إشعار، بريد) لتغيير تراجَع ولم يُحفَظ أصلًا. (هذا يفسّر ترتيب «احفظ ثمّ انشر» في معالِج
  قبول الحجز.)
\item \textbf{ترجمة التعارض}: لو ضرب الحفظ فهرسًا فريدًا (رقما الخطأ \en{2601}/\en{2627}
  في \en{SQL Server})، نلتقطه ونعيد رميه كـ \code{ConflictException} ودّيّ (مثل «هذا
  الموعد حُجِز للتوّ») بدل خطأ 500 غامض.
\end{itemize}

\begin{idea}[الفكرة: تحويل كل FK نحو المستخدم إلى Restrict]
الضمانة الرابعة تعيش في \code{OnModelCreating} لا في الحفظ، لكنّها من نفس فلسفة «قاعدة
مركزيّة واحدة». يرفض \en{SQL Server} وجود أكثر من مسار \en{CASCADE} إلى الجدول نفسه —
وهو ما يحدث لأيّ كيان له مفتاحان أجنبيّان أو أكثر نحو \code{ApplicationUser}. فبدل إضافة
\code{Restrict} في كل إعداد على حدة، نمرّ على \emph{كل} مفتاح أجنبيّ يشير للمستخدم ونحوّله
إلى \code{DeleteBehavior.Restrict} دفعةً واحدة — فنتجنّب الخطأ ونحمي تاريخ المستخدم.
\end{idea}

\begin{secnote}
تجميع هذه القواعد في نقطة واحدة أمانٌ بحدّ ذاته: لا يمكن لمطوّر أن «ينسى» ختم تاريخ أو
يترك تعارضًا يتسرّب كخطأ 500 يكشف تفاصيل داخليّة. القاعدة تُفرَض على الجميع لأنّها تعيش في
الطريق الذي لا مفرّ منه.
\end{secnote}

% ====================================================================
\subsection{ضوابط أخرى: التزامن والتكرار والفهارس الفريدة}

يبقى ثلاثة حرّاس صغار لكن حاسمين:

\subsubsection{التزامن المتفائل عبر RowVersion}
الكيانات القابلة للتدقيق تحمل عمودًا \code{RowVersion} يتغيّر مع كل تعديل. لو حاول
مستخدمان تعديل الصفّ نفسه في آنٍ، ينجح الأوّل ويفشل الثاني لأنّ نسخته صارت قديمة — فنتجنّب
«التحديث الضائع» دون أقفال ثقيلة. يُعرَّف بسطر:
\code{b.Property(p => p.RowVersion).IsRowVersion();}

\subsubsection{مفاتيح عدم التكرار IdempotencyKey}
المال والإشعارات يجب أن يحدثا \emph{مرّةً واحدة بالضبط} حتى مع إعادة المحاولة. لذلك يحمل
\code{Payment} و\code{Notification} حقل \code{IdempotencyKey} فريدًا. لو أُعيد إرسال نفس
الطلب، يصطدم بالمفتاح الموجود فلا يتكرّر خصمٌ ولا إشعار. (تذكّر \code{idempotencyKey} في
مقتطف سحب الدفع سابقًا — نفس الفكرة تُرسَل إلى \en{Stripe} كذلك.)

\subsubsection{الفهارس الفريدة المُصفّاة}
هذه أذكاها. أحيانًا نريد قاعدة مثل «تطبيق \emph{فعّال} واحد لكل طالب لكل منحة» — لكن نسمح
بعدّة تطبيقات مسحوبة أو مرفوضة. الفهرس الفريد العاديّ لا يكفي، فنستخدم \emph{فهرسًا فريدًا
مُصفّى}: يُطبّق شرط التفرّد على صفوف بعينها فقط.

\textbf{هنشوف إيه:} الفهرس الذي يفرض تطبيقًا فعّالًا واحدًا لكل طالب/منحة.

\begin{english}
\begin{lstlisting}[language={[Sharp]C}]
b.HasIndex(a => new { a.StudentId, a.ScholarshipId })
    .IsUnique()
    .HasFilter(
        "[ScholarshipId] IS NOT NULL " +
        "AND [Status] <> 'Withdrawn' " +
        "AND [Status] <> 'Rejected' " +
        "AND [Status] <> 'Accepted'")
    .HasDatabaseName("UX_Applications_Student_Scholarship_Active");
\end{lstlisting}
\end{english}

الشرط في \code{HasFilter} يستثني الحالات المنتهية (مسحوب/مرفوض/مقبول) والتطبيقات الخارجيّة
البحتة (بلا رابط منحة). فيبقى التفرّد مفروضًا على التطبيقات \emph{النشطة} فقط. النتيجة:
قاعدة العمل صارت محفورةً في القاعدة نفسها، لا مجرّد فحص في الكود يمكن تجاوزه.

\begin{qa}
\textbf{سؤال متوقّع:} «لو الكود فيه فحص قبل الإدراج، فلماذا نحتاج الفهرس الفريد أيضًا؟»
لأنّ الفحص في الكود عرضة لسباق التزامن (\en{race condition}): طلبان متزامنان قد يمرّان
كلاهما من الفحص ثم يُدرِجان. الفهرس الفريد في القاعدة هو خطّ الدفاع الأخير الذي لا
يُخترَق: أحد الطلبين يفشل حتمًا، ونحن نترجم فشله إلى \code{ConflictException} ودّيّ كما
رأينا في المعترِض. الكود يعطي رسالة لطيفة، والقاعدة تعطي الضمان.
\end{qa}

% ====================================================================
\begin{summary}[الخلاصة في سطور]
\begin{itemize}
\item \textbf{\en{CQRS}}: كل عملية إمّا \en{Command} (كتابة) أو \en{Query} (قراءة)،
  تُرسَل عبر \en{MediatR}، والـ \en{controller} رفيع بسطر واحد \code{mediator.Send}.
\item \textbf{المعالِج} هو موطن منطق العمل: رأينا \code{AcceptBookingCommand} يحمّل ثم
  يسحب الدفع ثم يجهّز غرفة \en{ACS} ثم يحفظ ثم ينشر حدث الدومين.
\item \textbf{خطّ المعالجة} أربع حلقات بترتيب مقصود: \en{Logging} ثم \en{Validation} ثم
  \en{Performance} ثم \en{Audit} (الأقرب للمعالِج) — كلٌّ في موضعه لسبب.
\item \textbf{التدقيق} تلقائيّ بسطر \code{[Auditable]} فوق الأمر؛ حلقة \en{Audit} تكتب صفّ
  \code{AuditLog} بعد نجاح المعالِج فقط.
\item \textbf{الحذف الناعم} عبر \code{ISoftDeletable} و\code{HasQueryFilter} على قرابة 15
  كيانًا يجعل الصفوف المحذوفة تختفي من كل استعلام دون تدميرها.
\item \textbf{معترِض \code{SaveChangesAsync}} يختم التواريخ، وينشر الأحداث بعد الحفظ فقط،
  ويترجم تعارض الفهرس الفريد إلى \code{ConflictException}، ويحوّل مفاتيح المستخدم إلى
  \code{Restrict}.
\item \textbf{ضوابط أخرى}: \code{RowVersion} للتزامن، \code{IdempotencyKey} لعدم التكرار،
  وفهارس فريدة مُصفّاة تحفر قواعد العمل داخل القاعدة نفسها.
\end{itemize}
\end{summary}
